% I. Предметна област, разгледани решения и обосновка на избора.
\section{Предметна област и обосновка на избора}
\label{sec:literature}

\subsection{Теоретична рамка}

Системата за компютърно зрение е верига от източник на данни, предварителна обработка,
извличане на признаци и вземане на решение. Апаратната част определя дефектите на входа:
пасивната стереодвойка дава плътна карта на дълбочината само там, където сцената има
текстура, а активните сензори дават плътност, независима от текстурата, но внасят шум по
краищата и загуба на данни при отразяващи повърхности. Информационната структура следва от
този избор: подреденото изображение на дълбочината запазва съседството по индекс, докато
облакът от точки го губи и налага пространствен индекс.

Предварителната обработка цели потискане на шума без разрушаване на геометрията. Линейното
изглаждане размива именно ръбовете, които следващите стъпки търсят, затова се предпочитат
методи, съобразени с ръба: двустранното филтриране претегля съседите едновременно по
разстояние и по разлика в стойността~\cite{tomasi1998bilateral}. Откриването на контури
следва постановката на Кани с трите ѝ критерия за добро откриване, добра локализация и
единичен отговор~\cite{canny1986edge}; върху карта на дълбочината същият апарат откроява
геометричните прекъсвания, но изисква отделно третиране на липсващите стойности.

Извличането на признаци разделя подредените от неподредените данни. За облаци от точки
утвърдено средство са хистограмите на бързите точкови признаци, които описват околността чрез
разпределение на ъглите между нормалите на съседите~\cite{rusu2009fpfh}. Сегментацията се води
или от геометрията, чрез съгласуваност на нормалите и кривината, или от подгонване на
параметричен модел, което почти винаги стъпва върху устойчиво оценяване със случайни
извадки~\cite{fischler1981ransac}. Текстурната сегментация ползва статистиките от матрицата на
съвместната поява на нива на сивото~\cite{haralick1973texture}.

Възстановяването на дълбочината по стереодвойка се свежда до задача за съответствие;
полуглобалното съгласуване я решава чрез сумиране на едномерни динамични програми по няколко
направления~\cite{hirschmuller2008sgm}. Когато сцената се снема от няколко положения, облаците
трябва да се приведат в обща координатна система. Каноничното решение е итеративната
най-близка точка~\cite{besl1992icp}, чиито варианти се различават по избора на съответствия,
по метриката и по стратегията на отхвърляне, и тези решения влияят на скоростта на сходимост
повече от самата формулировка~\cite{rusinkiewicz2001icp}. Всяка итерация изисква търсене на
най-близък съсед, което връща изложението към структурата на данните: k-d
дървото~\cite{bentley1975kdtree} остава основното средство, а приблизителните варианти купуват
скорост срещу контролирана загуба на точност~\cite{muja2014flann}. Разпознаването по модел
заема последното ниво и е извън реализацията; сегментираните области и описанията на
околностите са точно входът, който един разпознавател очаква. Утвърдените реализации с отворен
код~\cite{rusu2011pcl,bradski2000opencv} служат за ориентир кои алгоритми са установената
практика.

\subsection{Разгледани решения и направен избор}

\textbf{Пространствен индекс.} Вокселната решетка е тривиална за построяване и бърза при
равномерна плътност, но при неравномерна плътност или изразходва памет за празни кубчета,
или събира прекалено много точки в едно. K-d дървото~\cite{bentley1975kdtree} се приспособява
към плътността, защото разделя по медиана, но обхождането му е с нередовен достъп до паметта,
а октодървото стои между двете. Избрано е k-d дървото, а вокселната решетка се използва в
другата си роля, като стъпка за прореждане преди индексирането. Приблизителните
схеми~\cite{muja2014flann} са извън обхвата, защото внасят втори източник на грешка, който би
се смесил с грешката на измервания алгоритъм.

\textbf{Векторизация.} Вътрешните функции на конкретен набор инструкции дават пълен контрол,
но обвързват изходния текст с една архитектура. Библиотека за преносима векторизация запазва
единен изходен текст, но е външна зависимост. Автоматичната векторизация е най-евтина, но
зависи от компилатора. Избран е третият път, с добавка, която премахва главния му недостатък:
вътрешният цикъл е написан нарочно в пакетна форма, която компилаторът разпознава, и че я
разпознава, се проверява с \code{-fopt-info-vec} (Раздел~\ref{sec:implementation}). Условието
изборът да работи е координатите да са подредени по компонента, а не по точка, което превръща
оформлението на паметта в проектно решение (Раздел~\ref{sec:design}).

\textbf{Конвейер и външни библиотеки.} Монолитната функция е най-кратка, но прави отделната
стъпка неизмерима, затова всяка стъпка е обект с еднакъв договор, който приема облак и връща
облак заедно с показатели за изпълнението си. За външните библиотеки бяха разгледани две
положения: те да поемат входа и изхода и да служат за
еталон~\cite{rusu2011pcl,bradski2000opencv}, или проектът да няма нито една задължителна
външна зависимост. Избрано е второто, по практически довод: изграждане, което изисква
мениджър на пакети, на практика не се изпълнява от никого при първия опит. Цената е платена в
собствен код за четците, проверката, измерването и разлагането на матрица 3x3.
